\documentclass[12pt]{article}
\usepackage{subcaption}
\usepackage{graphicx} % Required for inserting images

\title{Lab 8: Neutron Measurements and Shielding}
\author{Owen Strong \\
\itshape Prof.: Angela DiFulvio\\
\itshape TA: Kholod Mahmoud \\
\itshape TA: Shaffer Bauer \\
\itshape TA: Justin Jia \\
}
\date{April 11, 2026}


\usepackage[
    style=ieee
]{biblatex}
\bibliography{refs.bib}

\usepackage[letterpaper, margin=1in]{geometry}
\usepackage{xspace}

% Figure Macros
\newcommand{\figwidth}{0.75\linewidth}
\newcommand{\figref}[1]{Fig.~\ref{#1}}
\newcommand{\defplotW}[3]{
    \begin{figure}[!htb]
        \centering
        \includegraphics[width=#3\linewidth]{fig/#1.png}
        \caption{#2}
        \newcommand{\tmplabel}{#1}
        \label{\tmplabel}
    \end{figure}
}
\newcommand{\defplot}[2]{\defplotW{#1}{#2}{0.75}}
\newcommand{\defsubplotW}[3]{
    % \hfill
    \begin{subfigure}{#3\linewidth}
        \includegraphics[width=\textwidth]{fig/#1.png}
        \caption{#2}
        \newcommand{\tmplabel}{#1}
        \label{\tmplabel}
    \end{subfigure}
    % \hfill
}
\newcommand{\defsubplot}[2]{\defsubplotW{#1}{#2}{0.45}}

% Common Macros
\newcommand{\micCi}{$\mathrm{\mu Ci}$\xspace}
\newcommand{\CsSrc}{$\mathrm{^{137}Cs}$\xspace}
\newcommand{\PuBe}{$\mathrm{^{238}PuBe}$\xspace}


\newcommand{\XSr}{\Sigma_R}
\newcommand{\incm}{$\mathrm{cm^{-1}}$\xspace}

\usepackage{placeins} % FloatBarrier
\usepackage{amsmath} % align
\usepackage{xcolor}
\usepackage[colorlinks=true, citecolor=violet, linkcolor=violet, urlcolor=violet]{hyperref}
\usepackage{enumitem}

\begin{document}
\pagenumbering{gobble}
\maketitle
\pagebreak
{\hypersetup{linkcolor=black}\tableofcontents}
\pagebreak
\pagenumbering{arabic}
\begin{abstract}
    Neutron radiation dose to radiation workers is a strong function of not only the sources of neutrons but the various shielding effects of materials in the immediate surroundings.
    In this laboratory, we characterize fast neutron shielding in common materials: concrete, steel, and polyethylene, and neutron reflection off of concrete shielding.
    We measure fast neutron dose rates using a MERIDIAN dose detector at a constant distance and measurement time attenuated through varying thicknesses of each material, then for the same thicknesses of concrete shielding with additional concrete shielding behind the source, behind both the source and the detector, and behind the detector alone.
    We measure macroscopic fast neutron removal cross-sections $\XSr$ of 0.047~\incm in concrete, 0.079~\incm in steel, and 0.108~\incm in polyethylene.
    Our observations most closely follow an exponential model for steel ($R^2=1.000$) and concrete ($R^2=0.999$), but deviate slightly for polyethylene ($R^2=0.975$), which match the expectations of neutron interaction for high-z materials like steel and low-z materials like polyethylene.
    We observe a strong relative dose with secondary shielding surrounding the detector or source without interstitial shielding  between the two (at a maximum of 1.284 times that without any shielding at all), and a noticeable relative dose with thick interstitial shielding with secondary shielding behind the detector, which we attribute to secondary reflection off of the environment.
    These results suggest the substantial contribution of neutron reflection off of concrete to nearby dose.
    Our results characterize fast neutron attenuation through concrete, steel, and polyethylene, and reflection off of concrete, and may be of interest to those designing shielding against high energy neutrons, and the safety and doses nearby concrete structures.
\end{abstract}

\section{Introduction and Background}

Neutron shielding is a critical phenomenon to designing safe and reliable nuclear devices.
The shielding of neutrons with different energies relates not only to the absorption neutrons at a given energy, but also the scattering of those neutrons to new, typically lower energies, resulting in an involved and nuanced shielding process.\\

A feasible simplification is to consider simply fast neutrons and slow, thermal neutrons, allowing this shielding process to be broken into three empirically-observed components: absorption of fast neutrons, scattering of fast neutrons into thermal neutrons, and the absorption of thermal neutrons.\\

The absorption of fast neutrons and scattering of fast neutrons to thermal neutrons might typically be combined into general ``removal'' of fast neutrons, and described using a semi-empirical formula, of exponential spacial attenuation:
\begin{equation}
    \label{eq:fast-remove}
    I(t) = I_0\exp(-\Sigma_rt)
\end{equation}
where $\Sigma_rt$ is a macroscopic cross-section for removal, and $I(t)$ the intensity of a neutron beam with initial intensity $I_0$, after attenuation through a thickness $t$.\\

\defplot{source_spectra}{Energy spectra of various common neutron sources. Figure sourced from lab manual.\cite{npre_lab_8_manual_2025}}

While for absorption of fast neutrons, the model of Eq.~\ref{eq:fast-remove} would be a strong fit over distance, the model might not be as high-fidelity for materials where downscattering interactions dominate the loss of fast neutrons.
As seen in \figref{source_spectra}, neutrons are often born at a wide range of energies, even before complex interaction with the environment, and for the same thickness of material different energies of neutrons might downscatter to different fast or slow energies, even before differences in cross-sections across energy are considered.\\

To directly measure $\XSr$, one could use a neutron detector with a high-energy threshold to measure only the flux or dose of fast neutrons.
However, in reality, such detectors are often inefficient, and more mathematically involved approaches (such as subsequently moderating the neutrons, and accounting for moderator absorption of fast neutrons\cite{npre_lab_8_manual_2025}) might instead be taken to make use of a thermal neutron detector.\\

In addition to these processes of attenuation through the material, another important consideration is the reflection of neutrons off of shielding.
An operator, standing behind the same shielding, might receive more dose if there is shielding close behind them, as neutrons are reflected off of the shield back in the operator's direction.\\

In this laboratory, we measure the attenuation of fast neutrons through varying thicknesses of iron, concrete, and polyethylene, in addition to the reflection of neutrons off of nearby concrete shielding, to characterize the removal cross section $\XSr$ of fast neutrons in these materials, and qualitatively characterize shielding and reflection more broadly.

\subsection{Dose Rate}
Consider a point source a sufficiently far distance $d$ from an aligned disc detector of radius $a$.
The geometric efficiency, or fractional solid angle, of this detector thus simplifies to $\Omega_f=\frac{a^2}{d^2}$\cite{knoll_radiation_2010}, allowing the dose $D(t)$ from a total source dose $D_0$ to be characterized using Eq.~\ref{eq:fast-remove}:
\begin{equation}
    \label{eq:disc-dose}
    D(t)=\left[D_0\frac{a^2}{d^2}\right]\exp(-\XSr t)
\end{equation}
allowing a linear relationship with slope $-\XSr$ to be observed for experimental $D(t)$, $d$, and $a$:
\begin{equation}
    \label{eq:log-dd}
    \begin{aligned}
        D(t)&=\left[D_0\frac{a^2}{d^2}\right]\exp(-\XSr t)\\
        \frac{d^2}{a^2}D(t)&=D_0\exp(-\XSr t)\\
        \ln\left[\frac{d^2}{a^2}D(t)\right]&=\ln(D_0) - \XSr t
    \end{aligned}
\end{equation}


\section{Experimental Procedure}\label{sec:methods}
In this experiment, we employed a MERIDIAN neutron detector\cite{noauthor_model_2003} to detect neutrons emitted by a \PuBe source. 
Specific equipment details may be found in Appendix~\ref{app:equip}.\\

\defplotW{Lab8_Setup}{The experimental setup used in this laboratory. For attenuation measurements, shielding combinations were placed in position~X. For reflection observations, concrete shielding blocks were additionally placed behind the source in position~Z, and in behind position~Y, as shown.}{0.5}

\defplot{Lab8_Diag}{A cross-sectional schematic of the experimental setup. Figure sourced from lab manual.\cite{npre_lab_8_manual_2025}}

We arranged the MERIDIAN detector and \PuBe source as shown in \figref{Lab8_Setup}, with additional supervision and care taken to safely handle the \PuBe source, such as by using the long attached rod to manipulate the source less directly.
In all measurements, the source was placed approximately 75~cm from the front of the detector, which given the length 24~cm of a MERIDIAN detector\cite{noauthor_model_2003}, yields an approximate distance $d=87$~cm.
Combined with a detector width of 22~cm, yielding radius $a=11$~cm, we estimate a geometric attenuation $\frac{a^2}{d^2}=0.0160$, which was held constant through all measurements.
During measurements when no interaction with the experimental setup was required, we ensured the half of the room adjacent to the source was unoccupied (i.e, no persons stood within approximately 5~meters of the source), to minimize unnecessary exposure to the \PuBe source.\\

We first measured attenuation of neutrons through different materials in position~X.
To perform each measurement, we arranged the desired thickness of material in a planar arrangement at this position, as seen in \figref{Lab8_Setup} for two blocks of concrete.
Resetting the counts on the MERIDIAN detector, setting acquisition time to 5~minutes and pressing the same button as to reset, we started a 5~minute integral count measurement.\\
% While waiting for the measurement to finish, we ensured no personnel stood within approximately 10 feet of the detector, to practice ALARA (As Low As Reasonable Achievable) radiation safety for the strong \PuBe source, and kept track of progress with a separate timer.\\

After completing this measurement, we quickly and carefully rearranged shielding materials to the next desired combination, and reset and repeated the measurement.

\subsection{Experiment 1: Material Comparison of Neutron Shielding}
For shielding comparison, we performed this radiation measurement procedure for 10~combinations of shielding: one measurement with no shielding, and measurements for shielding one, two and three layers thick each of concrete, steel, and polyethylene, placed in position~X.
The same bare measurement was used as a 0-thickness measurement for characterizations of all three materials, rather than recording a new measurement for each (lack of) material.\\

% For the layers of each material, we used concrete blocks of $10\pm0.1$~cm thickness each (for a maximum thickness of $30\pm0.3$~cm), pairs of steel plates $2\pm0.1$~cm thick each (up to total $6\pm0.3$~cm), and sets of polyethylene slabs $5\pm0.1$~cm thick (up to total $15\pm0.3$~cm). Reflective shielding placed in positions~Z and Y was composed of several blocks of the same $10\pm0.1$~cm concrete used at position~X.
For the layers of each material, we used concrete blocks of $\sim$10~cm thickness each (for a maximum thickness of $\sim$30~cm), pairs of steel plates $\sim$2~cm thick each (up to total $\sim$6~cm), and sets of polyethylene slabs $\sim$5~cm thick (up to total $\sim$15~cm). Reflective shielding placed in positions~Z and Y was composed of several blocks of the same $\sim$10~cm concrete used at position~X.

\subsection{Experiment 2: Secondary Shielding and Neutron Reflection}
For observation of neutron reflection, we used only concrete blocks, comparing combinations of concrete shielding in position~X (as in experiment 1) with combinations of shielding arranged behind the source in position~Z and position~Y (as seen in Figs.~\ref{Lab8_Setup} and \ref{Lab8_Diag}).
As in experiment 1, each series of measurements was taken with no shielding, and one, two, and three layers of concrete.\\

We first arranged a large shielding block behind the source in position~Z using around 10 blocks of concrete blocks, and performed measurements for the four thicknesses of shielding in position~X.\\

We then arranged a shielding block behind the detector (position~Y), keeping the position~Z shielding in place, to observe reflection both from shielding in Y and Z positions, and performed the four measurements again.\\

Finally, we removed the shielding block in Z to observe only reflection from Y, and performed the four measurements a final time.\\

To compare to concrete shielding at position~X with no reflection, the results of experiment~1, with identical thicknesses of concrete, were reused.

    
% \end{table}
\section{Experimental Results}\label{sec:results}
We successfully undertook radiation counts corresponding to each combination of shielding and reflection described above, and in this section detail such results, first as they relate to fast neutron attenuation, then as they relate to shield reflection behavior.

\subsection{Experiment 1: Material Comparison of Shielding}
\FloatBarrier

The doses received through each combination of shielding at position X are tabulated in Table~\ref{tab:atten-counts}, and logarithmically plotted (accounting for geometric attenuation $\frac{a^2}{d^2}$) in \figref{fig:xsr-plots}.
Using a linear regression defined in Eq.~\ref{eq:log-dd}, we estimate macroscopic removal cross-sections $\XSr$ for fast neutrons of 0.047~\incm in concrete (based on a correlation $R^2=0.999$), of 0.079~\incm in steel ($R^2=1.000$), and of 0.108~\incm in polyethylene ($R^2=0.975$).

\newcommand{\murem}{$\mathrm{\mu rem}$\xspace}
\begin{table}[!htb]
    \centering
    \caption{The integrated dose count, measured in \murem, by the MERIDIAN detector over a five-minute period with different combinations of shielding placed in position~X. The 0~Layer thick (no shielding) measurement, being independent of the shield material, was reused for all three materials.}
    \label{tab:atten-counts}
    \begin{tabular}{c|cccc}
        \hline
        Material & 0~Layers (cm) & 1~Layer & 2~Layer & 3~Layer \\
        \hline
        Concrete & 70.0 & 41.6 & 26.9 & 17.0 \\
        Steel & \textit{70.0} & 60.0 & 50.9 & 43.6 \\
        Polyethylene & \textit{70.0} & 47.9 & 29.2 & 13.6 \\
        \hline
    \end{tabular}
\end{table}

\begin{figure}
    \centering
    \defsubplotW{XSr_all}{For all materials compared.}{0.45}\hfill
    \defsubplot{XSr_conc}{For concrete.}\hfill
    \defsubplot{XSr_steel}{For steel.}\hfill
    \defsubplot{XSr_pe}{For polyethylene.}
    \caption{Plots of attenuation thickness $t$ vs $ln(\frac{d^2}{a^2}D(t))$ through each of the measured materials. In each plot, `$\Sigma R$' denotes the removal cross-section $\XSr$ estimated from measurements for that material.}
    \label{fig:xsr-plots}
\end{figure}

\FloatBarrier
\subsection{Experiment 2: Secondary Shielding and Neutron Reflection}
For each of the discussed combinations of concrete shielding in position~X and secondary concrete shielding in positions~Y and Z, integrated dose rates measured over the five-minute period are tabulated in 

\begin{table}[!htb]
    \centering\caption{The integrated dose count by the MERIDIAN detector over a five-minute period with different combinations of concrete shielding placed in positions~X (for thicknesses ranging from 0 to 30~cm), Y and Z (present or not present), relative to measurements with only shielding in position~X. Measurements of position~X shielding alone were reused from experiment~1's concrete tests, of identical configuration.}
    \label{tab:refl-counts}
    \begin{tabular}{c|cccc}
        \hline
        Shielding Positions & 0~Layers X & 1~Layer X & 2~Layers X & 3~Layers X \\\hline
        \it X (absolute) & \textit{70.0~$\mu rem$} & \textit{41.6~$\mu rem$} &\it 26.9~$\mu rem$ &\it 17.0~$\mu rem$ \\\hline
        X & 1 & 1 & 1 & 1 \\
        XZ & 1.123 & 1.014 & 1.063 & 1.047 \\
        XYZ & 1.284 & 0.974 & 0.967 & 1.141 \\
        XY & 1.126 & 0.990 & 1.052 & 1.147 \\\hline
        % XZ & 78.6 & 42.2 & 28.6 & 17.8 \\
        % XYZ & 89.9 & 40.5 & 26.0 & 19.4 \\
        % XY & 78.8 & 41.2 & 28.3 & 19.5 
    \end{tabular}
\end{table}
\FloatBarrier

\section{Discussion}\label{sec:disc}
All three materials seem to exhibit substantial $\XSr$ which could be feasibly estimated using Eq.~\ref{eq:log-dd} (0.047~\incm in concrete, 0.079~\incm in steel, and the largest, 0.108~\incm in polyethylene).
While the doses attenuated through concrete and steel shielding appear to strongly follow the behavior prescribed by this attenuation model (having linear correlation $R^2=0.999,1.000$ respectively, and meeting at the same intercept in \figref{XSr_all}), the observations of polyethylene have a weaker correlation $R^2=0.975$ (and do not meet at the same intercept).
This may relate to the neutron interactions expected in each medium.\\

As defined, $\XSr$ comprises both interactions where fast neutrons are absorbed, and those where they are scattered to lower energy.
The absorption of fast neutrons would be expected to very closely follow the attenuation model of Eq~\ref{eq:fast-remove}, but for scattering, as neutrons do not simply have discrete ``fast'' or ``slow'' energy, downscattered fast neutrons scattered in thin shielding might still have ``fast'' enough energy should they deflect towards the detector, while in thicker shielding the downscattering would then become more significant as neutrons begin to fall below the ``fast'' threshold.
Thus, if neutrons lose very little energy to scattering, $\XSr$ might remain sufficiently constant but if they lose much energy, an increasingly drastic might instead be observed.\\

This would suggest that in concrete and steel, fast neutron absorption dominates the effects seen in $\XSr$, which appears to remain constant, while in polyethylene, there is substantial contribution of neutron downscattering, causing $\XSr$ to become more substantial at higher thickness in \figref{XSr_pe}.
This is what we expect based on their compositions: concrete and steel are composed of high-z elements like Fe, Si, and Ca, for which neutrons lose less energy in scattering collisions\cite{yip_nuclear_2025}, while the organic polyethylene is composed of low-z elements like C and especially H, for which neutrons can lose the most energy in scattering\cite{yip_nuclear_2025}.\\

For the observations using secondary concrete shielding in positions~Y and Z, we observe a marked increase in dose for measurements without position~X shielding in Table~\ref{tab:refl-counts}.
This suggests that while it may not slow down (or ``moderate'') neutrons quickly enough to affect linear models of $\XSr$, neutron scattering is still a substantial phenomenon in concrete, with the highest relative dose (1.284) being measured with both Y and Z secondary shielding, but no X primary shielding.\\

As expected from such scattering behavior, the contributions of Y and Z secondary shielding are less strongly consistent to the log-linear attenuation, with relative doses being high for no X shielding, low for $\sim$10~cm of concrete X shielding, and increasing again for thicker X shielding (Table~\ref{tab:refl-counts}).
Contribution most noticeably increases again at $\sim$30~cm for the setups with added Y shielding behind the detector.\\

For Z, YZ, and Y shielding without X shielding, the scattered neutrons have a clear path to the detector, while with significant X shielding, the neutrons scattered off of Z shielding behind the source would still be attenuated substantially through the main X shielding, while any neutrons scattered off of X shielding would be unlikely to then scatter across the smaller shield of Y, and then again into the detector.
The increase in relative dose around $t\approx30$~cm with Y shielding behind the detector may be a result of scattering of neutrons first off of the nearby environment (such as the nearby metal chairs and cabinets in \figref{Lab8_Setup}), then off of Y shielding into the detector, which would not be affected by X shielding.
Additionally, while not expected to be as significant as for the lowest-z material polyethylene, some common atoms in concrete (like O, Al) are low enough that downscattering of neutrons to slow speeds (as opposed to absorption or scattering away from the detector) may still be significant contributions to phenomena with thicker concrete shielding.
\section{Conclusions}\label{sec:conc}

The objectives of this lab were largely met.
Neutron doses as attenuated through varying thicknesses of concrete, steel, and polyethylene were successfully measured, and used to estimate $\XSr$, and the effect of secondary concrete shielding reflecting behind the detector and source were successfully compared to intermediate concrete shielding.\\

Using a log-linear regression (thickness $t$ vs $\ln(D(t)\frac{d^2}{a^2})$) fast neutron $\XSr$ were measured of 0.047~\incm in concrete, 0.079~\incm in steel, and 0.108~\incm in polyethylene.
Dose rates had high correlation with the log-linear model for steel ($R^2=1.000$) and concrete ($R^2=0.999$), but a weaker correlation in polyethylene ($R^2=0.975$), as expected from the low-z of polyethylene atoms, for which neutrons tend to rapidly lose energy through scattering.\\

Neutron reflection off of secondary concrete shielding was most prominent without primary, central shielding, peaking at a relative dose of 1.284 with secondary shielding both behind the detector and source, to that observed with no shielding at all.
For moderate thicknesses of primary concrete shielding, the neutron reflection off of shielding did not seem to substantially contribute to the relative dose, while for the thickest primary shielding, there was more significant relative dose with shielding directly behind the detector (1.141 with shielding behind source and 1.147 without), likely from secondary reflection off of the environment unaffected by primary shielding.\\

While this laboratory characterized human equivalent dose, a repetition with neutron energy spectroscopy might yield further insight about the contributions of neutron scattering energy loss to $\XSr$.

% \bibliographystyle{sty/ieeetran}
\pagebreak\printbibliography

\appendix
\pagebreak
\FloatBarrier\section{Equipment Details}\label{app:equip}
In order to better replicate experiments and identify issues, the details of each equipment piece are recorded in Table~\ref{tab:equip}.

\newcommand{\na}{$N/A$}
\newcommand{\sref}[1]{$^{\ref{#1}}$}
\FloatBarrier
Further details of equipment are listed in Table~\ref{tab:equip}.
Where a detail could not be identified for a given device, it is replaced with ``\na''.
\begin{table}[!h]
    \centering
    \small
    \caption{Details of the equipment used in this laboratory.}
    \label{tab:equip}
    \scriptsize
    \begin{tabular}{c|cccc}
        Item & Manufacturer & Model & ID Type & ID Number\\\hline
        MERIDIAN Dose Detector & Far West Tech.\sref{addr:fwest} & Model 5085 & Property \# & PP10F79084 \\
        \PuBe Source & Unknown & Unknown & 1959 Activity & 480 mCi \\
        % Module rack & Canberra (now Mirion)\sref{addr:mirion}& Model 2000 & Inventory \# & NE759377 \\
        % Geiger-Muller Detector & \na & \na & \na & \na \\
        % GM Pulse Inverter & Spectrum Techniques\sref{addr:spect} & 2017 GPI & \na & \na \\
        % Multichannel Analyzer & Ortec\sref{addr:ortec} & EasyMCA & Inventory \# & P10F79090 \\
        % Alpha Spectrometer & Ortec\sref{addr:ortec} & Model 7401, Si surface barrier detector & Property \# & NE 759377\\
        % Alpha Source & Unknown & 0.1 \micCi \AmSrc & \na & \na\\
        % Vacuum Pump & Unknown & \na & \na & \na \\
        % Mylar Sheets & Unknown & $5\times3.6$~microns thick & \na & \na \\
        % Aluminum Sheets & Unkown & $5\times18$~microns thick & \na & \na \\
        % Weak Source 1 & Spectrum Techniques\sref{addr:spect} & 1.0\micCi ~$^{137}$Cs Beta Source & Date & Sep. 2017 \\
        % Weak Source 2 & Spectrum Techniques\sref{addr:spect} & 1.0\micCi ~$^{137}$Cs Beta Source & Date & Sep. 2017 \\
        % Strong Source & Eckert and Ziegler\sref{addr:eck-zieg} & 30\micCi~$^{137}Cs$ Beta Source & Date & Sep. 2017 \\
        % \InMeta~Source & \na & Irradiated Indium Foil & Date & Feb. 2026 \\
        % High Voltage Power Supply & Caen\sref{addr:caen} & N1470AL & Inventory \# & P10G96054 \\
        % Pulser & Ortec\sref{addr:ortec} & Model 480 & \na & \na \\
        % Oscilloscope & Keysight\sref{addr:keys} & Infiniivision DSO-X 2002A & Inventory \# & P10R03513 \\
        % Counter & Ortec\sref{addr:ortec} & Model 871 & Serial \# & 1024 \\
        % Preamplifier & Ortec\sref{addr:ortec} & Model 142PC & Serial \# & 2780r17 \\
        % Amplifier & Ortec\sref{addr:ortec} & Model 590A & Serial \# & 16076205 \\
        % Computer & Dell\sref{addr:dell} & Precision 360 & Eng. IT ID \# & EWS 20426 \\

    \end{tabular}
    \begin{enumerate}[label=\alph*]
        \item\label{addr:fwest}
        Far West Technology, Inc.
        (health physics instruments division):
        606 E Main,
        Puyallup, WA 98372;
        Phone: 805.964.3615;
        \url{https://www.fwt.com/index.htm}
        % \item\label{addr:mirion}
        % Mirion Technologies, Inc.:
        % 1218 Menlo Drive,
        % Atlanta, GA;
        % Phone: 770.432.2744;
        % \url{https://ir.mirion.com/}
        % \item\label{addr:spect}
        % Spectrum Techniques:
        % 106 Union Valley Rd, Oak Ridge, TN 37830;
        % Phone: 865.482.9937;
        % \url{https://www.spectrumtechniques.com/}
        % \item\label{addr:eck-zieg}
        % Eckert \& Ziegler Isotope Products:
        % 24937 Avenue Tibbitts, Valencia, CA 91355;
        % Phone: +1 866 476-9767;
        % \url{https://isotopeproducts.com/}
        % \item\label{addr:ortec}
        % Advanced Measurement Technology:
        % 801 South Illinois Avenue,
        % Oak Ridge, Tennessee 37830,
        % United States;
        % Phone: +1.865.482.4411;
        % Fax: +1.865.483.0396;
        % Email: ortec.info@ametek.com
        % \url{https://www.ortec-online.com};
        % \item\label{addr:caen}
        % Caen Sp.A.:
        % Via della Vetraia, 11, 55049 Viareggio LU - Italy;
        % Phone: +39 0584 388398;
        % \url{https://www.caen.it/}
        % \item\label{addr:keys}
        % Keysight Technologies: 
        % 1900 Garden of the Gods Road, Colorado Springs, CO 80907-3423 United States;
        % Phone: 1 800 829-4444;
        % \url{https://www.keysight.com}
        % \item\label{addr:dell}
        % Dell Technologies:
        % Dell Technologies
        % 1 Dell Way
        % Round Rock, TX 78664.;
        % Phone: 1-877-275-3355;
        % \url{https://www.dell.com/en-us/lp/contact-us}
    \end{enumerate}
\end{table}
\FloatBarrier
\end{document}
